\subsection{Why UFD Macro and OOD Pooled Rankings Differ}

LiteSSM-B obtains a higher \textbf{OOD Pooled} AUC (\textbf{0.722}) than LiteSSM-A (\textbf{0.712}), yet a lower \textbf{UFD Macro} (\textbf{0.700} vs.\ \textbf{0.718}).
Pooled AUC weights every OOD sample equally and can amplify easier or larger subsets inside \texttt{test\_ood.jsonl}.
Macro AUC averages generator-level AUCs with equal weight and remains sensitive to consistently weak generators such as DALL·E.
We therefore treat \textbf{UFD Macro as the primary cross-generator summary} for deployment-oriented reporting and keep pooled OOD as a secondary aggregate that must not substitute for Macro in the text or tables.

\subsection{SSM as a Stronger Deployment Point---Not a Universal Detector}

Across ID, UFD Macro, and bedroom-domain ID, LiteSSM-A occupies a favorable accuracy--size region under the evaluated latency budget.
LiteSSM-B is close on ID and competitive on several OOD views, but is heavier and slower; the LiteSSM-A--LiteSSM-B ID gap is CI-overlapping.
These gains do \emph{not} imply a universal or generator-agnostic detector: worst-generator AUC remains near chance for both SSMs.
The appropriate claim is narrower---under a matched lightweight protocol, SSM backbones improved average transfer relative to compact CNN baselines, while generator-specific collapse persisted.
LiteSSM-A is preferred for its highest Panel~A UFD Macro together with a locked batch-32 throughput of 226 images/s, not because it minimizes absolute batch-1 latency.

\subsection{Boundaries Exposed by Bedroom Content and DALL·E}

Near-ceiling CelebA-HQ AUCs show that aggregate ID scores can look strong even when a substantial content domain remains difficult (Section~\ref{sec:domain}).
Bedroom performance ($\approx$0.68--0.81) is where models separate and where SSM gains are most visible.
DALL·E exposes an orthogonal \textbf{generator} boundary: chance-level AUCs across CNN and SSM alike indicate that the limitation is not an artifact of LiteFreqNet alone (Section~\ref{sec:ufd}).

\subsection{Frequency Enhancement Is Not Universally Helpful}

LiteFreqNet~v2 did not win ID, UFD Macro, bedroom, or JPEG uniqueness.
Attaching a similar frequency branch to SSM backbones slightly hurt bake-off ID/OOD Pooled metrics.
Under the evaluated JPEG Q70 recompression setting, all compact models moved by at most 0.003 AUC (Section~\ref{sec:jpeg}); stability was shared rather than frequency-specific.
We therefore interpret frequency features as protocol- and architecture-dependent, not as a general recipe for lightweight AIGC detection.

\subsection{Limitations}

Training used a fixed no-pretraining recipe, a single training seed, and two DF training sources.
Bootstrap intervals quantify test-sample uncertainty but do not capture training stochasticity.
UFD Macro depends on the chosen subset list and per-class caps.
JPEG analysis focuses on Q70 for the main claim.
External reference detectors (UnivFD, NPR) use publicly released pretrained weights and native preprocessing; they contextualize performance and are not training-matched competitors.
An unmatched FLUX subset is reported only in Appendix~\ref{app:flux}~\cite{flux2024}.
Reported SSMs are study-specific pure-PyTorch classifiers for controlled comparison, not official architecture releases or vendor-kernel peaks.
Deepfake video and legal-forensics claims are out of scope.

\subsection{Practical Takeaway}

For practitioners selecting one compact still-image detector from the training-matched suite, \textbf{LiteSSM-A} is the recommended operating point under the evaluated latency budget, paired with explicit monitoring for difficult generators (e.g., DALL·E-like sources) and content domains (e.g., indoor scenes), rather than an assumption of uniform reliability.
